\honest{Moore's Paradox}{Eliezer Yudkowsky}{2009}

I think I understand Moore's paradox a bit better now, after reading the comments. I do not say what the paradox is. \nb{It is the puzzle of why a sentence like ``It is raining, but I do not believe that it is raining'' seems absurd to assert, though it is consistent and may be true. Moore's own example was ``I went to the pictures last Tuesday, but I don't believe that I did.''}

The commenter jimrandomh suggested that many people cannot tell ``I believe X'' from ``X.'' I disagree: young children understand what a false belief is. \nb{Research on false-belief tasks supports this.} But I offer a similar idea: many people ``may not consciously distinguish'' believing something from endorsing it. ``I believe in democracy'' means that you endorse democracy, so the word ``belief'' has more than one meaning, and a confused word may cause confused thinking.

So the woman from my earlier posts found reasons why it would be good to believe that people are nice, felt warmly toward that belief, and took the warmth for belief. Seeing that people are not so nice, she said ``I believe people are nicer than they are.'' I state this as what happened. \nb{Three days earlier the same sentence was explained by her belief that she had deceived herself. The post calls this one ``another mechanism.''} This is almost an honest mistake, since no one is taught how to tell when they believe something. The person who says the dragon in the garage is invisible does not notice that expecting to see no dragon means having a model with no dragon in it. What actual belief feels like is that a statement ``just seems like the way the world is''; that differs from feeling good about a belief you hold in quotation marks.

This makes ``this real-life case of Moore's Paradox'' seem less alien. \nb{If she meant ``endorse,'' her sentence is not a case of Moore's paradox as philosophers pose it: ``I endorse believing that people are nice, but they are not'' is neither inconsistent nor odd. The explanation shows how someone can utter a sentence of that form, not why the sentence is absurd when ``believe'' means believe.}

The commenter Kurige believes that God instilled a sense of right and wrong in us, and also that morality evolved. ``I suspect'' that Kurige has reasons to endorse both, and I ask whether they seem like good communities to join. \nb{Kurige's comment mentions neither communities nor benefits.} I ask Kurige to say both claims aloud without ``I believe'' and notice that they are harder to swallow. That difference is the difference between endorsing two beliefs and having one model of the world. \nb{Dropping ``I believe'' also drops a hedge, so the flat version may be harder for anyone to say. Kurige replied that saying ``God exists'' is harder for the same reason that a scientist would rather say ``according to our model, the higgs-boson should exist'' than ``the higgs-boson exists.''}
